\section{myco\_apps\_releases}\label{myco_apps_releases}

\href{https://www.r-project.org/}{\pandocbounded{\includesvg[keepaspectratio]{https://img.shields.io/badge/R-\%3E\%3D4.0-blue.svg}}}\\
\hyperref[-applications-disponibles]{\pandocbounded{\includesvg[keepaspectratio]{https://img.shields.io/badge/Domain-Mycologie\%20\%7C\%20\%C3\%89cologie-green.svg}}}\\
\hyperref[-applications-disponibles]{\pandocbounded{\includesvg[keepaspectratio]{https://img.shields.io/badge/Status-Op\%C3\%A9rationnel-brightgreen.svg}}}

Dépôt de publication des scripts R prêts à l'emploi issus des projets de
mycologie et d'écologie.\\
Chaque application est autonome : données d'exemple, script principal et
documentation inclus.

\textbf{Auteur :} Eddy Boite\\
\textbf{Dernière mise à jour :} 15 Juin 2026

\begin{center}\rule{0.5\linewidth}{0.5pt}\end{center}

\subsection{Table des matières}\label{table-des-matiuxe8res}

\begin{itemize}
\tightlist
\item
  \hyperref[-applications-disponibles]{Applications disponibles}
\item
  \hyperref[-structure-du-dux5cux25C3ux5cux25A9pux5cux25C3ux5cux25B4t]{Structure
  du dépôt}
\item
  \hyperref[-evaluation-du-potentiel-fongique-de-lintux5cux25C3ux5cux25A9rux5cux25C3ux5cux25AAt-patrimonial-et-du-gradient-chegd]{Évaluation
  du potentiel fongique, de l'intérêt patrimonial et du gradient CHEGD}

  \begin{itemize}
  \tightlist
  \item
    \hyperref[-quick-start-chegd]{Quick Start CHEGD}
  \item
    \hyperref[-structure-attendue-des-donnux5cux25C3ux5cux25A9es-chegd]{Structure
    attendue des données CHEGD}
  \item
    \hyperref[-lancer-le-script-chegd--guide-dux5cux25C3ux5cux25A9taillux5cux25C3ux5cux25A9]{Lancer
    le script CHEGD -- guide détaillé}
  \item
    \hyperref[-sorties-gux5cux25C3ux5cux25A9nux5cux25C3ux5cux25A9rux5cux25C3ux5cux25A9es-chegd]{Sorties
    générées CHEGD}
  \item
    \hyperref[-installation-des-dux5cux25C3ux5cux25A9pendances-chegd]{Installation
    des dépendances CHEGD}
  \item
    \hyperref[-configuration-chegd]{Configuration CHEGD}
  \item
    \hyperref[-interprux5cux25C3ux5cux25A9tation-des-indicateurs-chegd]{Interprétation
    des indicateurs CHEGD}
  \item
    \hyperref[-lecture-des-graphiques-chegd]{Lecture des graphiques
    CHEGD}
  \item
    \hyperref[-dux5cux25C3ux5cux25A9pannage-faq-chegd]{Dépannage (FAQ
    CHEGD)}
  \item
    \hyperref[-glossaire-chegd]{Glossaire CHEGD}
  \item
    \hyperref[-rux5cux25C3ux5cux25A9fux5cux25C3ux5cux25A9rences-chegd]{Références
    CHEGD}
  \end{itemize}
\item
  \hyperref[-inventaires-fongiques--complux5cux25C3ux5cux25A9tude--reprux5cux25C3ux5cux25A9sentativitux5cux25C3ux5cux25A9]{Inventaires
  fongiques -- Complétude \& Représentativité}

  \begin{itemize}
  \tightlist
  \item
    \hyperref[-prux5cux25C3ux5cux25A9requis-systux5cux25C3ux5cux25A8me-complets]{Prérequis
    système}
  \item
    \hyperref[-quick-start]{Quick Start}
  \item
    \hyperref[-vux5cux25C3ux5cux25A9rifier-que-tout-fonctionne]{Vérifier
    que tout fonctionne}
  \item
    \hyperref[-structure-attendue-des-donnux5cux25C3ux5cux25A9es]{Structure
    attendue des données}
  \item
    \hyperref[ux5cux25EFux5cux25B8ux5cux258F-lancer-le-script--guide-dux5cux25C3ux5cux25A9taillux5cux25C3ux5cux25A9]{Lancer
    le script -- guide détaillé}
  \item
    \hyperref[-sorties-gux5cux25C3ux5cux25A9nux5cux25C3ux5cux25A9rux5cux25C3ux5cux25A9es]{Sorties
    générées}
  \item
    \hyperref[-installation-des-dux5cux25C3ux5cux25A9pendances]{Installation
    des dépendances}
  \item
    \hyperref[-configuration]{Configuration}
  \item
    \hyperref[-interprux5cux25C3ux5cux25A9tation-des-indicateurs]{Interprétation
    des indicateurs}
  \item
    \hyperref[-lecture-des-graphiques]{Lecture des graphiques}
  \item
    \hyperref[-bonnes-pratiques-dinterprux5cux25C3ux5cux25A9tation]{Bonnes
    pratiques}
  \item
    \hyperref[-dux5cux25C3ux5cux25A9pannage-faq]{Dépannage (FAQ)}
  \item
    \hyperref[-historique-des-versions]{Historique des versions}
  \item
    \hyperref[-glossaire]{Glossaire}
  \item
    \hyperref[-Rux5cux25C3ux5cux25A9fux5cux25C3ux5cux25A9rences]{Références}
  \end{itemize}
\end{itemize}

\begin{center}\rule{0.5\linewidth}{0.5pt}\end{center}

\subsection{Applications disponibles}\label{applications-disponibles}

{\def\LTcaptype{none} % do not increment counter
\begin{longtable}[]{@{}
  >{\raggedright\arraybackslash}p{(\linewidth - 8\tabcolsep) * \real{0.2000}}
  >{\raggedright\arraybackslash}p{(\linewidth - 8\tabcolsep) * \real{0.2000}}
  >{\raggedright\arraybackslash}p{(\linewidth - 8\tabcolsep) * \real{0.2000}}
  >{\raggedright\arraybackslash}p{(\linewidth - 8\tabcolsep) * \real{0.2000}}
  >{\raggedright\arraybackslash}p{(\linewidth - 8\tabcolsep) * \real{0.2000}}@{}}
\toprule\noalign{}
\begin{minipage}[b]{\linewidth}\raggedright
\#
\end{minipage} & \begin{minipage}[b]{\linewidth}\raggedright
Application
\end{minipage} & \begin{minipage}[b]{\linewidth}\raggedright
Script
\end{minipage} & \begin{minipage}[b]{\linewidth}\raggedright
Version
\end{minipage} & \begin{minipage}[b]{\linewidth}\raggedright
Domaine
\end{minipage} \\
\midrule\noalign{}
\endhead
\bottomrule\noalign{}
\endlastfoot
1 &
\hyperref[-inventaires-fongiques--complux5cux25C3ux5cux25A9tude--reprux5cux25C3ux5cux25A9sentativitux5cux25C3ux5cux25A9]{Inventaires
fongiques -- Complétude \& Représentativité} &
\texttt{scripts/Inventaires\_completude\_representativite.R} & v1.4 &
Mycologie / Écologie \\
2 &
\hyperref[-evaluation-du-potentiel-fongique-de-lintux5cux25C3ux5cux25A9rux5cux25C3ux5cux25AAt-patrimonial-et-du-gradient-chegd]{Évaluation
du potentiel fongique, de l'intérêt patrimonial et du gradient CHEGD} &
\texttt{scripts/Evaluation\_Potentiel\_Fongique\_Interets\_Patrimoniaux\_CHEGD.R}
& v1.0 & Mycologie / Patrimonialité \\
\end{longtable}
}

\begin{center}\rule{0.5\linewidth}{0.5pt}\end{center}

\subsection{Structure du dépôt}\label{structure-du-duxe9puxf4t}

\begin{verbatim}
myco_apps_releases/
├── data/
│   └── observations.csv                  # Données d'exemple (prêtes à utiliser)
├── docs/
│   ├── Readme_Analyse_Inventaires_Fongiques.md            # Documentation détaillée (ICR)
│   ├── Inventaires_completude_representativite_Conception_Technique.md
│   ├── Evaluation_Potentiel_Fongique_Interets_Patrimoniaux_CHEGD_Conception_Fonctionnelle.tex
│   └── Evaluation_Potentiel_Fongique_Interets_Patrimoniaux_CHEGD_Conception_technique.tex
├── logs/                                 # Journaux d'exécution
├── results/
│   ├── ICR/                             # Dossier de sortie (créé automatiquement)
│   └── EPFIP_CHEGD/                     # Sorties de l'évaluation CHEGD
├── scripts/
│   ├── Inventaires_completude_representativite.R          # Script principal (v1.4)
│   └── Evaluation_Potentiel_Fongique_Interets_Patrimoniaux_CHEGD.R
└── README.md                             # Ce fichier
\end{verbatim}

\begin{center}\rule{0.5\linewidth}{0.5pt}\end{center}

\subsection{Évaluation du potentiel fongique, de l'intérêt patrimonial
et du gradient
CHEGD}\label{uxe9valuation-du-potentiel-fongique-de-lintuxe9ruxeat-patrimonial-et-du-gradient-chegd}

Pipeline R dédié à l'évaluation multi-critères des pelouses sur trois
axes :

\begin{itemize}
\tightlist
\item
  \textbf{Potentiel fongique} (score et classe),
\item
  \textbf{Intérêt patrimonial} (indice et classe),
\item
  \textbf{Gradient CHEGD} et \textbf{indice de représentativité (IR)}.
\end{itemize}

Le script intègre également un module de \textbf{fiabilité de
détermination} (objectifs 1 à 4), avec export des statistiques, figures
et synthèses de sélection de modèles.

\textbf{Script principal :}
\texttt{scripts/Evaluation\_Potentiel\_Fongique\_Interets\_Patrimoniaux\_CHEGD.R}

\textbf{Conceptions associées :}

\begin{itemize}
\tightlist
\item
  \texttt{docs/Evaluation\_Potentiel\_Fongique\_Interets\_Patrimoniaux\_CHEGD\_Conception\_Fonctionnelle.tex}
\item
  \texttt{docs/Evaluation\_Potentiel\_Fongique\_Interets\_Patrimoniaux\_CHEGD\_Conception\_technique.tex}
\end{itemize}

\subsubsection{Entrées attendues}\label{entruxe9es-attendues}

Le script supporte :

\begin{itemize}
\tightlist
\item
  \texttt{CSV} (séparateur auto-détecté)
\item
  \texttt{XLSX} (avec ou sans feuille explicitée)
\end{itemize}

Colonnes métier attendues (configurables dans
\texttt{get\_embedded\_config()}) :

\begin{itemize}
\tightlist
\item
  \texttt{Espèces}
\item
  \texttt{Famille}
\item
  \texttt{Date}
\item
  \texttt{Nombre\ d\textquotesingle{}espèce}
\item
  \texttt{Site}
\item
  \texttt{Fiabilité\ détermination}
\end{itemize}

Par défaut, le fichier d'entrée est :

\begin{itemize}
\tightlist
\item
  \texttt{data/données\_récoltes\_chegd\_pelouses.csv}
\end{itemize}

\subsubsection{Lancer l'analyse CHEGD}\label{lancer-lanalyse-chegd}

En ligne de commande, depuis la racine du dépôt :

\texttt{Rscript\ scripts/Evaluation\_Potentiel\_Fongique\_Interets\_Patrimoniaux\_CHEGD.R}

\subsubsection{Sorties principales}\label{sorties-principales}

Les résultats sont exportés dans :

\begin{itemize}
\tightlist
\item
  \texttt{results/EPFIP\_CHEGD/}
\end{itemize}

Fichiers de sortie majeurs :

\begin{itemize}
\tightlist
\item
  \texttt{resume\_global.csv}
\item
  \texttt{resume\_par\_site.csv}
\item
  \texttt{potentiel\_fongique\_par\_site.csv}
\item
  \texttt{indice\_patrimonial\_par\_site.csv}
\item
  \texttt{gradient\_chegd\_par\_site.csv}
\item
  \texttt{indice\_representativite\_ir\_par\_site.csv}
\item
  \texttt{synthese\_evaluation\_par\_site.csv}
\end{itemize}

Figures générées :

\begin{itemize}
\tightlist
\item
  \texttt{fig1\_tableau\_de\_bord\_pelouses} (PNG/PDF)
\item
  \texttt{fig2\_positionnement\_ecologique} (PNG/PDF)
\item
  \texttt{fig3\_composition\_potentiel} (PNG/PDF)
\item
  \texttt{fig4\_chegd\_par\_visite} (PNG/PDF)
\item
  \texttt{fig5\_heatmap\_classes} (PNG/PDF)
\item
  \texttt{fig6\_niveaux\_fiabilite} (PNG/PDF)
\item
  \texttt{fig7\_indice\_representativite\_ir} (PNG/PDF)
\end{itemize}

Le journal d'exécution est écrit dans :

\begin{itemize}
\tightlist
\item
  \texttt{logs/EPFIP\_CHEGD\_YYYYMMDD\_HHMM.log}
\end{itemize}

\begin{center}\rule{0.5\linewidth}{0.5pt}\end{center}

\subsubsection{Quick Start CHEGD}\label{quick-start-chegd}

\textbf{Pour lancer rapidement l'application CHEGD avec les données du
dépôt :}

Cette section suit le même format pas-à-pas que l'application
Inventaires, afin de faciliter la prise en main des deux pipelines.

\paragraph{Étape 1️⃣ --- Vérifier le fichier
d'entrée}\label{uxe9tape-1-vuxe9rifier-le-fichier-dentruxe9e}

Par défaut, le script lit :

\begin{itemize}
\tightlist
\item
  \texttt{data/données\_récoltes\_chegd\_pelouses.csv}
\end{itemize}

Vous pouvez conserver ce nom, ou modifier la configuration embarquée
dans le script.

\paragraph{Étape 2️⃣ --- Lancer le
script}\label{uxe9tape-2-lancer-le-script}

Depuis la racine du dépôt :

\texttt{Rscript\ scripts/Evaluation\_Potentiel\_Fongique\_Interets\_Patrimoniaux\_CHEGD.R}

\paragraph{Étape 3️⃣ --- Consulter les
sorties}\label{uxe9tape-3-consulter-les-sorties}

Les résultats sont produits dans :

\begin{itemize}
\tightlist
\item
  \texttt{results/EPFIP\_CHEGD/}
\end{itemize}

avec les tableaux CSV, les figures PNG/PDF et le journal d'exécution
dans \texttt{logs/}.

\begin{center}\rule{0.5\linewidth}{0.5pt}\end{center}

\subsubsection{Structure attendue des données
CHEGD}\label{structure-attendue-des-donnuxe9es-chegd}

Le script accepte \textbf{CSV} ou \textbf{XLSX} (séparateur CSV
auto-détecté).

\paragraph{Colonnes métier
attendues}\label{colonnes-muxe9tier-attendues}

{\def\LTcaptype{none} % do not increment counter
\begin{longtable}[]{@{}
  >{\raggedright\arraybackslash}p{(\linewidth - 4\tabcolsep) * \real{0.3333}}
  >{\raggedright\arraybackslash}p{(\linewidth - 4\tabcolsep) * \real{0.3333}}
  >{\raggedright\arraybackslash}p{(\linewidth - 4\tabcolsep) * \real{0.3333}}@{}}
\toprule\noalign{}
\begin{minipage}[b]{\linewidth}\raggedright
Colonne
\end{minipage} & \begin{minipage}[b]{\linewidth}\raggedright
Type attendu
\end{minipage} & \begin{minipage}[b]{\linewidth}\raggedright
Description
\end{minipage} \\
\midrule\noalign{}
\endhead
\bottomrule\noalign{}
\endlastfoot
\texttt{Espèces} & Texte & Nom de l'espèce observée \\
\texttt{Famille} & Texte & Famille taxonomique \\
\texttt{Date} & Date ou texte convertible & Date d'observation \\
\texttt{Nombre\ d\textquotesingle{}espèce} & Numérique & Abondance /
effectif observé \\
\texttt{Site} & Texte & Libellé du site (ex. \texttt{Pelouse\ 16}) \\
\texttt{Fiabilité\ détermination} & Texte & Niveau de confiance de
détermination \\
\end{longtable}
}

\begin{quote}
Ces colonnes sont configurables dans \texttt{get\_embedded\_config()}.
\end{quote}

\paragraph{Formats de date reconnus}\label{formats-de-date-reconnus}

Le script gère notamment :

\begin{itemize}
\tightlist
\item
  \texttt{YYYY-MM-DD}
\item
  \texttt{DD/MM/YYYY}
\item
  \texttt{DD-MM-YYYY}
\item
  dates Excel (numériques)
\end{itemize}

\begin{center}\rule{0.5\linewidth}{0.5pt}\end{center}

\subsubsection{Lancer le script CHEGD -- guide
détaillé}\label{lancer-le-script-chegd-guide-duxe9tailluxe9}

Méthodes d'exécution identiques à la section Inventaires : une option en
ligne de commande pour la reproductibilité et une option interactive
pour l'exploration.

\paragraph{Option A --- Ligne de commande
(recommandée)}\label{option-a-ligne-de-commande-recommanduxe9e}

\begin{verbatim}
cd /chemin/vers/myco_apps_releases
Rscript scripts/Evaluation_Potentiel_Fongique_Interets_Patrimoniaux_CHEGD.R
\end{verbatim}

\paragraph{Option B --- RStudio / session
interactive}\label{option-b-rstudio-session-interactive}

\begin{verbatim}
setwd("/chemin/vers/myco_apps_releases")
source("scripts/Evaluation_Potentiel_Fongique_Interets_Patrimoniaux_CHEGD.R")
\end{verbatim}

\paragraph{Ce qui se passe pendant
l'exécution}\label{ce-qui-se-passe-pendant-lexuxe9cution}

Le script :

\begin{enumerate}
\def\labelenumi{\arabic{enumi}.}
\tightlist
\item
  lit et valide les données d'entrée,
\item
  calcule les indicateurs par site (potentiel, patrimonialité, CHEGD),
\item
  calcule les indicateurs de fiabilité (objectifs 1 à 4),
\item
  exporte CSV + figures,
\item
  écrit un log horodaté dans \texttt{logs/}.
\end{enumerate}

Si un classeur de référence est détecté (feuilles dédiées), le pipeline
active un \textbf{mode d'alignement ``fidélité Excel''} pour reproduire
les valeurs métier du classeur.

\begin{center}\rule{0.5\linewidth}{0.5pt}\end{center}

\subsubsection{Sorties générées
CHEGD}\label{sorties-guxe9nuxe9ruxe9es-chegd}

Les sorties sont écrites dans :

\begin{itemize}
\tightlist
\item
  \texttt{results/EPFIP\_CHEGD/}
\end{itemize}

\paragraph{Fichiers de synthèse
principaux}\label{fichiers-de-synthuxe8se-principaux}

{\def\LTcaptype{none} % do not increment counter
\begin{longtable}[]{@{}
  >{\raggedright\arraybackslash}p{(\linewidth - 2\tabcolsep) * \real{0.5000}}
  >{\raggedright\arraybackslash}p{(\linewidth - 2\tabcolsep) * \real{0.5000}}@{}}
\toprule\noalign{}
\begin{minipage}[b]{\linewidth}\raggedright
Fichier
\end{minipage} & \begin{minipage}[b]{\linewidth}\raggedright
Description
\end{minipage} \\
\midrule\noalign{}
\endhead
\bottomrule\noalign{}
\endlastfoot
\texttt{resume\_global.csv} & Vue d'ensemble du jeu analysé \\
\texttt{resume\_par\_site.csv} & Statistiques agrégées par site \\
\texttt{potentiel\_fongique\_par\_site.csv} & Score/classe de potentiel
fongique \\
\texttt{indice\_patrimonial\_par\_site.csv} & Indice/classe
patrimoniale \\
\texttt{gradient\_chegd\_par\_site.csv} & Gradient CHEGD par site et
visites \\
\texttt{indice\_representativite\_ir\_par\_site.csv} & Indice IR par
visite/site \\
\texttt{synthese\_evaluation\_par\_site.csv} & Consolidation
multi-indicateurs \\
\end{longtable}
}

\paragraph{Module fiabilité (Objectifs 1 à
4)}\label{module-fiabilituxe9-objectifs-1-uxe0-4}

Le script exporte aussi des fichiers statistiques dédiés à la fiabilité
de détermination, notamment :

\begin{itemize}
\tightlist
\item
  distributions de niveaux de fiabilité,
\item
  comparaison de schémas de pondération,
\item
  sélection de modèles (validation croisée),
\item
  synthèse finale de recommandation.
\end{itemize}

\paragraph{Figures générées}\label{figures-guxe9nuxe9ruxe9es}

\begin{itemize}
\tightlist
\item
  \texttt{fig1\_tableau\_de\_bord\_pelouses}
\item
  \texttt{fig2\_positionnement\_ecologique}
\item
  \texttt{fig3\_composition\_potentiel}
\item
  \texttt{fig4\_chegd\_par\_visite}
\item
  \texttt{fig5\_heatmap\_classes}
\item
  \texttt{fig6\_niveaux\_fiabilite}
\item
  \texttt{fig7\_indice\_representativite\_ir}
\end{itemize}

Chaque figure est exportée en \textbf{PNG} et \textbf{PDF}.

\begin{center}\rule{0.5\linewidth}{0.5pt}\end{center}

\subsubsection{Installation des dépendances
CHEGD}\label{installation-des-duxe9pendances-chegd}

\paragraph{Packages requis}\label{packages-requis}

{\def\LTcaptype{none} % do not increment counter
\begin{longtable}[]{@{}ll@{}}
\toprule\noalign{}
Package & Rôle principal \\
\midrule\noalign{}
\endhead
\bottomrule\noalign{}
\endlastfoot
\texttt{readxl} & Lecture des fichiers Excel \\
\texttt{dplyr} & Manipulation de données \\
\texttt{stringr} & Traitement de chaînes \\
\texttt{ggplot2} & Visualisation \\
\texttt{gridExtra} & Composition multi-graphiques \\
\texttt{MASS} & Modèles statistiques ordonnés \\
\texttt{nnet} & Modèles multinomiaux \\
\end{longtable}
}

Commande groupée :

\begin{verbatim}
install.packages(c("readxl", "dplyr", "stringr", "ggplot2", "gridExtra", "MASS", "nnet"))
\end{verbatim}

\paragraph{Package optionnel}\label{package-optionnel}

\begin{itemize}
\tightlist
\item
  \texttt{ggrepel} (améliore le placement des labels sur certaines
  figures)
\end{itemize}

\begin{quote}
Le script tente d'installer automatiquement les packages requis s'ils
sont absents.
\end{quote}

\begin{center}\rule{0.5\linewidth}{0.5pt}\end{center}

\subsubsection{Configuration CHEGD}\label{configuration-chegd}

La configuration est intégrée dans la fonction
\texttt{get\_embedded\_config()} du script.

Paramètres principaux :

{\def\LTcaptype{none} % do not increment counter
\begin{longtable}[]{@{}
  >{\raggedright\arraybackslash}p{(\linewidth - 4\tabcolsep) * \real{0.3333}}
  >{\raggedright\arraybackslash}p{(\linewidth - 4\tabcolsep) * \real{0.3333}}
  >{\raggedright\arraybackslash}p{(\linewidth - 4\tabcolsep) * \real{0.3333}}@{}}
\toprule\noalign{}
\begin{minipage}[b]{\linewidth}\raggedright
Paramètre
\end{minipage} & \begin{minipage}[b]{\linewidth}\raggedright
Valeur par défaut
\end{minipage} & \begin{minipage}[b]{\linewidth}\raggedright
Description
\end{minipage} \\
\midrule\noalign{}
\endhead
\bottomrule\noalign{}
\endlastfoot
\texttt{input\_file} &
\texttt{data/données\_récoltes\_chegd\_pelouses.csv} & Fichier
d'entrée \\
\texttt{input\_sheet} & \texttt{NULL} & Feuille Excel à lire (si
XLSX) \\
\texttt{output\_dir} & \texttt{results} & Dossier parent de sortie \\
\texttt{output\_prefix} & \texttt{EPFIP\_CHEGD} & Sous-dossier de
sortie \\
\texttt{columns\$species} & \texttt{Espèces} & Colonne espèce \\
\texttt{columns\$family} & \texttt{Famille} & Colonne famille \\
\texttt{columns\$date} & \texttt{Date} & Colonne date \\
\texttt{columns\$count} & \texttt{Nombre\ d\textquotesingle{}espèce} &
Colonne abondance \\
\texttt{columns\$site} & \texttt{Site} & Colonne site \\
\texttt{columns\$reliability} & \texttt{Fiabilité\ détermination} &
Colonne fiabilité \\
\end{longtable}
}

Exemple d'adaptation :

\begin{verbatim}
# Dans get_embedded_config()
input_file = "data/mon_fichier_chegd.xlsx"
input_sheet = "Feuil1"
\end{verbatim}

\begin{center}\rule{0.5\linewidth}{0.5pt}\end{center}

\subsubsection{Interprétation des indicateurs
CHEGD}\label{interpruxe9tation-des-indicateurs-chegd}

\paragraph{Potentiel fongique}\label{potentiel-fongique}

Le score est classé selon trois niveaux :

\begin{itemize}
\tightlist
\item
  \textbf{≤ 10} : \texttt{Potentiel\ fongique\ faible}
\item
  \textbf{{]}10 ; 30{[}} : \texttt{Potentiel\ fongique\ intéressant}
\item
  \textbf{≥ 30} : \texttt{Potentiel\ fongique\ élevé}
\end{itemize}

\paragraph{Intérêt patrimonial}\label{intuxe9ruxeat-patrimonial}

Classement selon l'indice patrimonial :

\begin{itemize}
\tightlist
\item
  \textbf{≤ 2} : \texttt{Intérêt\ faible}
\item
  \textbf{3 à 5} : \texttt{Intérêt\ local}
\item
  \textbf{6 à 10} : \texttt{Intérêt\ régional}
\item
  \textbf{11 à 14} : \texttt{Intérêt\ national}
\item
  \textbf{≥ 15} : \texttt{Intérêt\ international}
\end{itemize}

\paragraph{Gradient CHEGD et IR}\label{gradient-chegd-et-ir}

Le script calcule un gradient par visite et un indice de
représentativité :

\[IR_{visite} = \max\left(0, 1 - \frac{gradient_{visite}}{chegd_{total}}\right)\]

\begin{itemize}
\tightlist
\item
  \textbf{IR proche de 1} : meilleure représentativité,
\item
  \textbf{IR faible} : marge de progression plus importante sur la
  couverture.
\end{itemize}

\begin{center}\rule{0.5\linewidth}{0.5pt}\end{center}

\subsubsection{Lecture des graphiques
CHEGD}\label{lecture-des-graphiques-chegd}

{\def\LTcaptype{none} % do not increment counter
\begin{longtable}[]{@{}
  >{\raggedright\arraybackslash}p{(\linewidth - 2\tabcolsep) * \real{0.5000}}
  >{\raggedright\arraybackslash}p{(\linewidth - 2\tabcolsep) * \real{0.5000}}@{}}
\toprule\noalign{}
\begin{minipage}[b]{\linewidth}\raggedright
Figure
\end{minipage} & \begin{minipage}[b]{\linewidth}\raggedright
Lecture principale
\end{minipage} \\
\midrule\noalign{}
\endhead
\bottomrule\noalign{}
\endlastfoot
\texttt{fig1\_tableau\_de\_bord\_pelouses} & Comparaison globale des 3
axes par pelouse \\
\texttt{fig2\_positionnement\_ecologique} & Positionnement potentiel ×
patrimonialité, couleur/taille = CHEGD \\
\texttt{fig3\_composition\_potentiel} & Décomposition du score potentiel
par groupes fonctionnels \\
\texttt{fig4\_chegd\_par\_visite} & Variation du gradient CHEGD selon
les visites \\
\texttt{fig5\_heatmap\_classes} & Vue décisionnelle synthétique des
classes finales \\
\texttt{fig6\_niveaux\_fiabilite} & Répartition des niveaux de fiabilité
de détermination \\
\texttt{fig7\_indice\_representativite\_ir} & Heatmap de l'IR par visite
et par pelouse \\
\end{longtable}
}

\begin{center}\rule{0.5\linewidth}{0.5pt}\end{center}

\subsubsection{Dépannage (FAQ CHEGD)}\label{duxe9pannage-faq-chegd}

\textbf{Le script ne trouve pas mon fichier d'entrée}\\
Vérifier \texttt{input\_file} dans \texttt{get\_embedded\_config()} et
la présence réelle du fichier dans \texttt{data/}.

\textbf{Erreur de colonnes manquantes}\\
Contrôler les intitulés exacts des colonnes métier (ou adapter
\texttt{columns\$...} dans la configuration).

\textbf{Dates mal interprétées}\\
Vérifier le format de la colonne \texttt{Date} et éviter les valeurs
hétérogènes dans une même colonne.

\textbf{Pourquoi mes valeurs diffèrent du classeur de référence ?}\\
Le mode ``fidélité Excel'' n'est activé que si un classeur de référence
compatible est détecté (feuilles attendues).

\textbf{Je n'ai pas toutes les figures}\\
Certaines figures dépendent de la disponibilité de données suffisantes
ou de certaines structures (ex. gradients par visite exploitables).

\begin{center}\rule{0.5\linewidth}{0.5pt}\end{center}

\subsubsection{Glossaire CHEGD}\label{glossaire-chegd}

{\def\LTcaptype{none} % do not increment counter
\begin{longtable}[]{@{}
  >{\raggedright\arraybackslash}p{(\linewidth - 2\tabcolsep) * \real{0.5000}}
  >{\raggedright\arraybackslash}p{(\linewidth - 2\tabcolsep) * \real{0.5000}}@{}}
\toprule\noalign{}
\begin{minipage}[b]{\linewidth}\raggedright
Terme
\end{minipage} & \begin{minipage}[b]{\linewidth}\raggedright
Définition
\end{minipage} \\
\midrule\noalign{}
\endhead
\bottomrule\noalign{}
\endlastfoot
\textbf{CHEGD} & Gradient composite basé sur les groupes fonctionnels
utilisés pour l'évaluation écologique des pelouses. \\
\textbf{IR} & Indice de représentativité calculé par visite :
\(IR = \max(0, 1 - gradient_{visite}/chegd_{total})\). \\
\textbf{Potentiel fongique} & Score synthétique basé sur la présence de
groupes et d'espèces indicatrices, puis classé en 3 niveaux. \\
\textbf{Indice patrimonial} & Indice de valeur patrimoniale du site,
classé de \emph{faible} à \emph{international}. \\
\textbf{Fiabilité détermination} & Niveau de confiance taxonomique de
l'observation (ex. certaine/probable/non renseignée). \\
\textbf{Mode fidélité Excel} & Mécanisme d'alignement des sorties script
sur un classeur de référence quand les feuilles attendues sont
détectées. \\
\end{longtable}
}

\begin{center}\rule{0.5\linewidth}{0.5pt}\end{center}

\subsubsection{Références CHEGD}\label{ruxe9fuxe9rences-chegd}

\begin{itemize}
\tightlist
\item
  Sellier, Y. et coll. --- \emph{Bulletin de la Société Mycologique de
  France}, vol.~131, pp.~1--2.
\item
  Protocole métier interne documenté dans :
\item
  \texttt{docs/Evaluation\_Potentiel\_Fongique\_Interets\_Patrimoniaux\_CHEGD\_Conception\_Fonctionnelle.tex}
\item
  \texttt{docs/Evaluation\_Potentiel\_Fongique\_Interets\_Patrimoniaux\_CHEGD\_Conception\_technique.tex}
\item
  Implémentation opérationnelle :
  \texttt{scripts/Evaluation\_Potentiel\_Fongique\_Interets\_Patrimoniaux\_CHEGD.R}.
\end{itemize}

\begin{center}\rule{0.5\linewidth}{0.5pt}\end{center}

\subsection{Inventaires fongiques -- Complétude \&
Représentativité}\label{inventaires-fongiques-compluxe9tude-repruxe9sentativituxe9}

Pipeline R pour automatiser l'analyse de la \textbf{complétude
d'inventaires fongiques} et de leur \textbf{représentativité}, via :

\begin{itemize}
\tightlist
\item
  le \textbf{TEE} (Taux d'Espèces Exceptionnelles)
\item
  l'\textbf{Ir} (Indice de représentativité : \(Ir = 1 - TEE\))
\item
  la courbe temps-espèces avec ajustement linéaire et hyperbolique
\item
  les métriques de pertinence scientifique (stabilité, robustesse,
  cohérence spatio-temporelle)
\item
  les comparaisons multi-sites avec classement et heatmap
\end{itemize}

Le script produit automatiquement des tableaux CSV et des graphiques PNG
(export PDF et SVG également disponibles), par site et en synthèse
globale, avec un manifeste automatique de traçabilité en fin
d'exécution.

Validation CSV renforcée incluse : contrôle de schéma, audit
\texttt{readr::problems()}, rapport de conformité exporté, et arrêt en
mode strict si non-conformité.

\begin{quote}
Documentation complète :
\url{docs/Readme_Analyse_Inventaires_Fongiques.md}\\
Conception technique :
\url{docs/Inventaires_completude_representativite_Conception_Technique.md}
\end{quote}

\begin{center}\rule{0.5\linewidth}{0.5pt}\end{center}

\subsubsection{Prérequis système
complets}\label{pruxe9requis-systuxe8me-complets}

\paragraph{Installer R (obligatoire)}\label{installer-r-obligatoire}

\textbf{R} est un environnement de calcul statistique gratuit et
open-source. Le script s'exécute avec \textbf{R ≥ 4.0} sur
\textbf{Windows}, \textbf{macOS} et \textbf{Linux}.

\subparagraph{Télécharger et installer
R}\label{tuxe9luxe9charger-et-installer-r}

\begin{enumerate}
\def\labelenumi{\arabic{enumi}.}
\tightlist
\item
  \textbf{Accéder au site officiel :}
  \href{https://cran.r-project.org}{cran.r-project.org}
\item
  \textbf{Sélectionner votre système d'exploitation :}

  \begin{itemize}
  \tightlist
  \item
    \textbf{Windows} : cliquer sur
    \href{https://cran.r-project.org/bin/windows/}{Download R for
    Windows}, puis
    \href{https://cran.r-project.org/bin/windows/base/}{base}, puis
    télécharger le dernier exécutable \texttt{.exe}
  \item
    \textbf{macOS} : cliquer sur
    \href{https://cran.r-project.org/bin/macosx/}{Download R for macOS},
    puis télécharger le \texttt{.pkg} correspondant à votre architecture
    (Intel ou Apple Silicon)
  \item
    \textbf{Linux} : suivre les instructions pour Debian/Ubuntu, Fedora,
    Red Hat, CentOS, etc.
  \end{itemize}
\item
  \textbf{Exécuter l'installateur} et suivre les instructions par défaut
  (tout en \textbf{Next} convient)
\item
  \textbf{Vérifier l'installation :} ouvrir un \textbf{Terminal}
  (macOS/Linux) ou \textbf{Invite de Commandes} (Windows) et taper :
\end{enumerate}

\begin{verbatim}
R --version
\end{verbatim}

Vous devez voir \texttt{R\ version\ 4.x.x} ou plus récent.

\paragraph{Installer RStudio Desktop (fortement recommandé, mais
optionnel)}\label{installer-rstudio-desktop-fortement-recommanduxe9-mais-optionnel}

\textbf{RStudio} est un \textbf{IDE} (interface graphique) qui rend R
beaucoup plus convivial pour les débutants. Ce n'est pas obligatoire,
mais \textbf{très utile}.

\subparagraph{Télécharger et installer
RStudio}\label{tuxe9luxe9charger-et-installer-rstudio}

\begin{enumerate}
\def\labelenumi{\arabic{enumi}.}
\tightlist
\item
  \textbf{Accéder au site :}
  \href{https://posit.co/download/rstudio-desktop/}{posit.co/download/rstudio-desktop}
\item
  \textbf{Télécharger la version gratuite} (version ``Community'',
  open-source)
\item
  \textbf{Sélectionner votre système d'exploitation} et télécharger
  l'installateur
\item
  \textbf{Exécuter l'installateur} et suivre les instructions par défaut
\item
  \textbf{Ouvrir RStudio} pour la première fois --- vous verrez une
  fenêtre avec 4 panneaux (Editeur, Console, Environnement, Fichiers)
\end{enumerate}

\paragraph{Installer les packages R
requis}\label{installer-les-packages-r-requis}

Les \textbf{packages} sont des extensions qui ajoutent des
fonctionnalités à R. Le script utilise une dizaine de packages pour
manipuler les données et générer les graphiques.

\subparagraph{Méthode A --- Via RStudio (recommandée pour les
débutants)}\label{muxe9thode-a-via-rstudio-recommanduxe9e-pour-les-duxe9butants}

\begin{enumerate}
\def\labelenumi{\arabic{enumi}.}
\tightlist
\item
  \textbf{Ouvrir RStudio}
\item
  \textbf{Dans le panneau en bas-à-droite} (onglet ``Packages''),
  cliquer sur le bouton \textbf{Install}
\item
  \textbf{Copier-coller} la liste de packages ci-dessous dans la zone
  qui s'ouvre :
\end{enumerate}

\begin{verbatim}
dplyr tidyr ggplot2 purrr readr stringr forcats tibble scales gridExtra
\end{verbatim}

\begin{enumerate}
\def\labelenumi{\arabic{enumi}.}
\tightlist
\item
  \textbf{Vérifier} que ``Install from: CRAN repository'' est
  sélectionné
\item
  \textbf{Cliquer sur ``Install''} et attendre (1--5 minutes selon la
  connexion)
\end{enumerate}

\subparagraph{Méthode B --- Via la console R
(texte)}\label{muxe9thode-b-via-la-console-r-texte}

\begin{enumerate}
\def\labelenumi{\arabic{enumi}.}
\tightlist
\item
  \textbf{Ouvrir la console R} (application ``R'' sur Windows/macOS, ou
  console classique)
\item
  \textbf{Copier-coller} cette commande complète et appuyer sur
  \textbf{Entrée} :
\end{enumerate}

\begin{verbatim}
install.packages(c(
  "dplyr", "tidyr", "ggplot2", "purrr",
  "readr", "stringr", "forcats", "tibble",
  "scales", "gridExtra"
))
\end{verbatim}

\begin{enumerate}
\def\labelenumi{\arabic{enumi}.}
\tightlist
\item
  \textbf{Répondre} \texttt{**y**} \textbf{ou} \texttt{**yes**} aux
  questions qui s'affichent
\item
  \textbf{Attendre} que l'installation se termine (aucune erreur rouge
  au final ne doit présager un problème ; une installation réussie se
  termine par \texttt{Done})
\end{enumerate}

\subparagraph{Vérifier l'installation}\label{vuxe9rifier-linstallation}

Après l'installation, exécuter cette commande pour vérifier que tout les
packages requis sont présents :

\begin{verbatim}
required_pkgs <- c("dplyr", "tidyr", "ggplot2", "purrr", "readr", "stringr", "forcats", "tibble", "scales", "gridExtra")
missing <- !sapply(required_pkgs, requireNamespace, quietly = TRUE)
if (any(missing)) {
  cat("❌ Packages manquants :", paste(required_pkgs[missing], collapse = ", \n"))
} else {
  cat("✅ Tous les packages requis sont installés !")
}
\end{verbatim}

\begin{center}\rule{0.5\linewidth}{0.5pt}\end{center}

\paragraph{Résumé : les 3 étapes pour se
préparer}\label{ruxe9sumuxe9-les-3-uxe9tapes-pour-se-pruxe9parer}

{\def\LTcaptype{none} % do not increment counter
\begin{longtable}[]{@{}
  >{\raggedright\arraybackslash}p{(\linewidth - 6\tabcolsep) * \real{0.2500}}
  >{\raggedright\arraybackslash}p{(\linewidth - 6\tabcolsep) * \real{0.2500}}
  >{\raggedright\arraybackslash}p{(\linewidth - 6\tabcolsep) * \real{0.2500}}
  >{\raggedright\arraybackslash}p{(\linewidth - 6\tabcolsep) * \real{0.2500}}@{}}
\toprule\noalign{}
\begin{minipage}[b]{\linewidth}\raggedright
\#
\end{minipage} & \begin{minipage}[b]{\linewidth}\raggedright
Étape
\end{minipage} & \begin{minipage}[b]{\linewidth}\raggedright
Durée
\end{minipage} & \begin{minipage}[b]{\linewidth}\raggedright
Difficulté
\end{minipage} \\
\midrule\noalign{}
\endhead
\bottomrule\noalign{}
\endlastfoot
\textbf{1} & Télécharger et installer \textbf{R} & 5--10 min & Facile
(suivre les instructions) \\
\textbf{2} & Télécharger et installer \textbf{RStudio} & 5--10 min &
Facile (optionnel mais recommandé) \\
\textbf{3} & Installer les \textbf{packages R} & 2--5 min & Facile
(copier-coller une commande) \\
\end{longtable}
}

\textbf{Résultat :} vous avez tout ce qu'il faut pour lancer le script !

\begin{center}\rule{0.5\linewidth}{0.5pt}\end{center}

\paragraph{Besoin d'aide ? Dépannage
initial}\label{besoin-daide-duxe9pannage-initial}

\textbf{``R n'est pas reconnu'' dans le Terminal/Invite de Commandes}

\begin{itemize}
\tightlist
\item
  \textbf{Windows} : redémarrer le Terminal \textbf{après}
  l'installation de R
\item
  \textbf{macOS/Linux} : la commande \texttt{R\ -\/-version} fonctionne
  généralement
\end{itemize}

\textbf{``je ne sais pas ouvrir un Terminal''}

\begin{itemize}
\tightlist
\item
  \textbf{Windows} : appuyer sur \textbf{Windows + R}, taper
  \texttt{cmd}, appuyer sur \textbf{Entrée}
\item
  \textbf{macOS} : \textbf{Cmd + Espace}, taper \texttt{Terminal},
  appuyer sur \textbf{Entrée}
\item
  \textbf{Linux} : \textbf{Ctrl + Alt + T} (dépend de la distribution)
\end{itemize}

\textbf{``Un package refuse de s'installer''}

\begin{itemize}
\tightlist
\item
  Ignorer pour l'instant --- le script affichera un message explicite si
  un package obligatoire manque
\item
  Réessayer plus tard avec une bonne connexion Internet
\end{itemize}

\textbf{``je veux lancer le script directement sans RStudio''}

\begin{itemize}
\tightlist
\item
  C'est possible, voir section
  \hyperref[ux5cux25EFux5cux25B8ux5cux258F-lancer-le-script--guide-dux5cux25C3ux5cux25A9taillux5cux25C3ux5cux25A9]{Lancer
  le script -- guide détaillé}, Option A (ligne de commande)
\end{itemize}

\begin{center}\rule{0.5\linewidth}{0.5pt}\end{center}

\subsubsection{Quick Start}\label{quick-start}

\textbf{Pour les débutants qui viennent d'installer R et RStudio :}

\paragraph{Étape 1️⃣ --- Préparer vos
données}\label{uxe9tape-1-pruxe9parer-vos-donnuxe9es}

Placer votre fichier d'observations dans le dossier \texttt{data/} du
dépôt, sous le nom \texttt{observations.csv}.

Un fichier d'exemple (\texttt{data/observations.csv}) est \textbf{déjà
fourni} --- vous pouvez l'utiliser directement pour tester !

\paragraph{Étape 2️⃣ --- Ouvrir RStudio et charger le
script}\label{uxe9tape-2-ouvrir-rstudio-et-charger-le-script}

\begin{enumerate}
\def\labelenumi{\arabic{enumi}.}
\tightlist
\item
  \textbf{Ouvrir RStudio}
\item
  \textbf{Cliquer sur File → Open File} et naviguer vers :
\item
  \textbf{Cliquer sur ``Source''} (bouton en haut-droite du panneau
  éditeur) pour lancer le script
\item
  \textbf{Attendre} que l'exécution se termine (vous verrez les messages
  de progression dans la console en bas)
\end{enumerate}

\paragraph{Étape 3️⃣ --- Consulter les
résultats}\label{uxe9tape-3-consulter-les-ruxe9sultats}

Les sorties seront créées dans le dossier \texttt{**results/**} :

\begin{itemize}
\tightlist
\item
  Fichiers \textbf{CSV} = tableaux de données à ouvrir avec Excel ou un
  tableur
\item
  Fichiers \textbf{PNG} = graphiques prêts à consulter
\end{itemize}

C'est tout ! 🎉

\paragraph{Alternative : ligne de commande (pour utilisateurs
avancés)}\label{alternative-ligne-de-commande-pour-utilisateurs-avancuxe9s}

Si vous êtes à l'aise avec le Terminal, ouvrir un terminal, naviguer
jusqu'à la racine du dépôt et exécuter :

\begin{verbatim}
Rscript scripts/Inventaires_completude_representativite.R
\end{verbatim}

\begin{center}\rule{0.5\linewidth}{0.5pt}\end{center}

\subsubsection{Vérifier que tout
fonctionne}\label{vuxe9rifier-que-tout-fonctionne}

Avant de lancer le script sur vos vraies données, un simple test de
vérification :

\paragraph{Test 1 : R fonctionne}\label{test-1-r-fonctionne}

Ouvrir \textbf{RStudio} (ou la console R), et copier-coller ceci dans la
console :

\begin{verbatim}
print("Hello from R! ✅")
\end{verbatim}

Si vous voyez \texttt{{[}1{]}\ "Hello\ from\ R!\ ✅"}, c'est bon.

\paragraph{Test 2 : Les packages sont bien
installés}\label{test-2-les-packages-sont-bien-installuxe9s}

Copier-coller ceci dans la console :

\begin{verbatim}
required_pkgs <- c("dplyr", "tidyr", "ggplot2", "purrr", "readr", "stringr", "forcats", "tibble", "scales", "gridExtra")
all_ok <- all(sapply(required_pkgs, requireNamespace, quietly = TRUE))
if (all_ok) {
  print("✅ Tous les packages obligatoires sont prêts !")
} else {
  print("❌ Certains packages manquent. Relancer l'installation.")
}
\end{verbatim}

Si vous voyez \texttt{✅}, vous êtes prêt à lancer le script.

\paragraph{Test 3 : Lancer le script sur les données
d'exemple}\label{test-3-lancer-le-script-sur-les-donnuxe9es-dexemple}

Maintenant, lancez le script en cliquant sur ``Source'' (comme dans le
Quick Start). Vous devez voir des messages horodatés dans la console,
puis un \textbf{manifeste} en fin d'exécution listant tous les fichiers
créés.

Si tout s'est bien passé, consultez le dossier \texttt{results/} ---
vous y trouverez les fichiers CSV et PNG.

\textbf{Félicitations ! Vous êtes prêt à lancer le script sur vos
propres données.}

\begin{center}\rule{0.5\linewidth}{0.5pt}\end{center}

\begin{center}\rule{0.5\linewidth}{0.5pt}\end{center}

\subsubsection{Structure attendue des
données}\label{structure-attendue-des-donnuxe9es}

Le fichier d'entrée doit être en format \textbf{CSV} (ou TSV / TXT,
détection automatique du séparateur).

\paragraph{Colonnes obligatoires}\label{colonnes-obligatoires}

{\def\LTcaptype{none} % do not increment counter
\begin{longtable}[]{@{}lll@{}}
\toprule\noalign{}
Colonne & Type & Description \\
\midrule\noalign{}
\endhead
\bottomrule\noalign{}
\endlastfoot
\texttt{site} & Texte & Identifiant du site inventorié \\
\texttt{date} & Date (\texttt{YYYY-MM-DD}) & Date de la visite \\
\texttt{visite\_id} & Texte & Identifiant de la visite (ex. \texttt{V1},
\texttt{V2}\ldots) \\
\texttt{espece} & Texte & Nom de l'espèce observée \\
\end{longtable}
}

\paragraph{Colonne optionnelle}\label{colonne-optionnelle}

{\def\LTcaptype{none} % do not increment counter
\begin{longtable}[]{@{}
  >{\raggedright\arraybackslash}p{(\linewidth - 4\tabcolsep) * \real{0.3333}}
  >{\raggedright\arraybackslash}p{(\linewidth - 4\tabcolsep) * \real{0.3333}}
  >{\raggedright\arraybackslash}p{(\linewidth - 4\tabcolsep) * \real{0.3333}}@{}}
\toprule\noalign{}
\begin{minipage}[b]{\linewidth}\raggedright
Colonne
\end{minipage} & \begin{minipage}[b]{\linewidth}\raggedright
Type
\end{minipage} & \begin{minipage}[b]{\linewidth}\raggedright
Description
\end{minipage} \\
\midrule\noalign{}
\endhead
\bottomrule\noalign{}
\endlastfoot
\texttt{placette} & Texte & Identifiant de la placette (active l'analyse
spatiale et la CA/AFC) \\
\end{longtable}
}

\begin{quote}
\textbf{Important :} une visite est identifiée de façon unique par la
combinaison \texttt{site\ +\ date\ +\ visite\_id}. Cela permet de gérer
les jeux de données où un même \texttt{visite\_id} peut être réutilisé à
des dates différentes.
\end{quote}

\paragraph{Exemple de données}\label{exemple-de-donnuxe9es}

{\def\LTcaptype{none} % do not increment counter
\begin{longtable}[]{@{}lllll@{}}
\toprule\noalign{}
site & date & visite\_id & espece & placette \\
\midrule\noalign{}
\endhead
\bottomrule\noalign{}
\endlastfoot
BoisLarge & 2025-09-10 & V1 & Gyrodon lividus & P1 \\
BoisLarge & 2025-09-10 & V1 & Russula leprosa & P1 \\
BoisLarge & 2025-09-17 & V2 & Amanita rubescens & P2 \\
Site\_B & 2025-10-01 & V1 & Cantharellus cibarius & P3 \\
\end{longtable}
}

Le fichier \texttt{data/observations.csv} fourni contient un jeu de
données d'exemple fonctionnel avec plusieurs sites et placettes.

\begin{center}\rule{0.5\linewidth}{0.5pt}\end{center}

\subsubsection{Lancer le script -- guide
détaillé}\label{lancer-le-script-guide-duxe9tailluxe9}

\paragraph{Option A --- Ligne de commande
(Rscript)}\label{option-a-ligne-de-commande-rscript}

Méthode recommandée pour une exécution reproductible et automatisable.

\begin{verbatim}
# Ouvrir un terminal, se positionner à la racine du dépôt
cd /chemin/vers/myco_apps_releases

# Lancer le script
Rscript scripts/Inventaires_completude_representativite.R
\end{verbatim}

\paragraph{Option B --- Session R interactive (RStudio ou console
R)}\label{option-b-session-r-interactive-rstudio-ou-console-r}

\begin{verbatim}
# Définir le répertoire de travail sur la racine du dépôt
setwd("/chemin/vers/myco_apps_releases")

# Charger et exécuter le script
source("scripts/Inventaires_completude_representativite.R")
\end{verbatim}

\begin{quote}
Dans RStudio, vous pouvez également ouvrir le fichier \texttt{.R} et
utiliser le bouton \textbf{Source} (Ctrl+Shift+S / Cmd+Shift+S).
\end{quote}

\paragraph{Contrôle de l'exécution via variable
d'environnement}\label{contruxf4le-de-lexuxe9cution-via-variable-denvironnement}

{\def\LTcaptype{none} % do not increment counter
\begin{longtable}[]{@{}
  >{\raggedright\arraybackslash}p{(\linewidth - 4\tabcolsep) * \real{0.3333}}
  >{\raggedright\arraybackslash}p{(\linewidth - 4\tabcolsep) * \real{0.3333}}
  >{\raggedright\arraybackslash}p{(\linewidth - 4\tabcolsep) * \real{0.3333}}@{}}
\toprule\noalign{}
\begin{minipage}[b]{\linewidth}\raggedright
Variable
\end{minipage} & \begin{minipage}[b]{\linewidth}\raggedright
Valeur
\end{minipage} & \begin{minipage}[b]{\linewidth}\raggedright
Effet
\end{minipage} \\
\midrule\noalign{}
\endhead
\bottomrule\noalign{}
\endlastfoot
\texttt{INVENTAIRES\_AUTO\_RUN} & \texttt{TRUE} (défaut) & Lance
l'analyse automatiquement au chargement \\
\texttt{INVENTAIRES\_AUTO\_RUN} & \texttt{FALSE} (ou \texttt{0},
\texttt{no}, \texttt{off}) & Charge uniquement les fonctions sans les
exécuter \\
\end{longtable}
}

\begin{verbatim}
# Charger les fonctions sans lancer l'analyse (utile pour tester ou déboguer)
INVENTAIRES_AUTO_RUN=FALSE Rscript scripts/Inventaires_completude_representativite.R
\end{verbatim}

\paragraph{Ce qui s'affiche dans la console pendant
l'exécution}\label{ce-qui-saffiche-dans-la-console-pendant-lexuxe9cution}

Le script journalise chaque étape avec horodatage
\texttt{{[}YYYY-MM-DD\ HH:MM:SS{]}\ {[}INFO{]}} :

\begin{verbatim}
[2026-03-23 10:00:01] [INFO] Initialisation du script ...
[2026-03-23 10:00:01] [INFO] Lecture du fichier : data/observations.csv
[2026-03-23 10:00:01] [INFO] Données valides : 120 lignes, 3 sites, 15 espèces
[2026-03-23 10:00:02] [INFO] Traitement site : BoisLarge (8 visites)
[2026-03-23 10:00:03] [INFO] Traitement site : Site_B (6 visites)
...
[2026-03-23 10:00:08] [INFO] ================ MANIFESTE DES SORTIES ================
[2026-03-23 10:00:08] [INFO] [MANIFEST] OK   results/ICR_donnees_preparees.csv
[2026-03-23 10:00:08] [INFO] [MANIFEST] OK   results/ICR_resume_tous_sites.csv
...
\end{verbatim}

En fin d'exécution, le \textbf{manifeste} liste l'état de chaque fichier
attendu :

\begin{itemize}
\tightlist
\item
  \texttt{OK} --- fichier produit avec succès
\item
  \texttt{MANQUANT} --- fichier obligatoire non généré (voir les logs
  pour l'erreur)
\item
  \texttt{OPTIONNEL\_NON\_GENERE} --- fichier optionnel non généré
  (conditions non réunies)
\end{itemize}

\paragraph{Options de débogage (à modifier directement en tête du
script)}\label{options-de-duxe9bogage-uxe0-modifier-directement-en-tuxeate-du-script}

{\def\LTcaptype{none} % do not increment counter
\begin{longtable}[]{@{}
  >{\raggedright\arraybackslash}p{(\linewidth - 4\tabcolsep) * \real{0.3333}}
  >{\raggedright\arraybackslash}p{(\linewidth - 4\tabcolsep) * \real{0.3333}}
  >{\raggedright\arraybackslash}p{(\linewidth - 4\tabcolsep) * \real{0.3333}}@{}}
\toprule\noalign{}
\begin{minipage}[b]{\linewidth}\raggedright
Option
\end{minipage} & \begin{minipage}[b]{\linewidth}\raggedright
Défaut
\end{minipage} & \begin{minipage}[b]{\linewidth}\raggedright
Effet si \texttt{TRUE}
\end{minipage} \\
\midrule\noalign{}
\endhead
\bottomrule\noalign{}
\endlastfoot
\texttt{DEBUG\_MODE} & \texttt{FALSE} & Affiche des messages très
détaillés et tous les warnings R \\
\texttt{BENCHMARK\_MODE} & \texttt{FALSE} & Mesure et affiche le temps
d'exécution de chaque calcul \\
\texttt{CLEAN\_ENVIRONMENT} & \texttt{FALSE} & Nettoie l'environnement R
avant exécution (usage isolé uniquement) \\
\end{longtable}
}

\begin{center}\rule{0.5\linewidth}{0.5pt}\end{center}

\subsubsection{Sorties générées}\label{sorties-guxe9nuxe9ruxe9es}

Toutes les sorties sont créées dans \texttt{results/ICR/} (le dossier
est créé automatiquement s'il n'existe pas).

\paragraph{Vue d'ensemble du pipeline (9 étapes →
fichiers)}\label{vue-densemble-du-pipeline-9-uxe9tapes-fichiers}

\begin{verbatim}
1) Préparation/validation des données
    -> Global CSV : ICR_donnees_preparees.csv
2) Richesse observée et cumulée (courbe temps-espèces)
    -> Site CSV   : <site>/ICR_01_courbe_temps_especes.csv
    -> Site PNG   : <site>/ICR_01_richesse_par_visite_et_cumul.png
3) Ajustements linéaire/hyperbolique (si conditions remplies)
    -> Site CSV   : <site>/ICR_05_modeles.csv
    -> Site PNG   : <site>/ICR_02_courbe_temps_especes_hyperbole.png
4) Asymptote hyperbolique et ratio de complétude
    -> Site CSV   : <site>/ICR_04_resume_site.csv
    -> Global CSV : ICR_resume_tous_sites.csv
    -> Global PNG : ICR_comparaison_completude_sites.png
5) TEE (taux d'espèces exceptionnelles) et indice Ir
    -> Site CSV   : <site>/ICR_02_tee_ir.csv
    -> Site PNG   : <site>/ICR_03_tee_ir.png
    -> Global PNG : ICR_comparaison_ir_sites.png
6) Fréquences d'occurrence des espèces
    -> Site CSV   : <site>/ICR_03_frequence_especes.csv
    -> Site PNG   : <site>/ICR_04_histogramme_frequences.png
7) Occupation spatiale (si placette disponible)
    -> Site CSV   : <site>/ICR_06_occupation_spatiale.csv
    -> Site PNG   : <site>/ICR_05_temporel_vs_spatial.png
8) CA/AFC placettes-espèces (si vegan et données adaptées)
    -> Site PNG   : <site>/ICR_06_CA_placettes_especes.png
9) Métriques de pertinence scientifique (stabilité, robustesse, cohérence)
    -> Site CSV   : <site>/ICR_07_metrics_pertinence.csv
    -> Global CSV : ICR_metrics_pertinence_tous_sites.csv
    -> Site PNG   : <site>/ICR_07_metrics_pertinence_dashboard.png
    -> Global PNG : ICR_metrics_pertinence_heatmap_sites.png
    -> Global PNG : ICR_metrics_pertinence_score_sites.png
\end{verbatim}

\paragraph{\texorpdfstring{Fichiers globaux
(\texttt{results/})}{Fichiers globaux (results/)}}\label{fichiers-globaux-results}

{\def\LTcaptype{none} % do not increment counter
\begin{longtable}[]{@{}
  >{\raggedright\arraybackslash}p{(\linewidth - 4\tabcolsep) * \real{0.3333}}
  >{\raggedright\arraybackslash}p{(\linewidth - 4\tabcolsep) * \real{0.3333}}
  >{\raggedright\arraybackslash}p{(\linewidth - 4\tabcolsep) * \real{0.3333}}@{}}
\toprule\noalign{}
\begin{minipage}[b]{\linewidth}\raggedright
Fichier
\end{minipage} & \begin{minipage}[b]{\linewidth}\raggedright
Type
\end{minipage} & \begin{minipage}[b]{\linewidth}\raggedright
Description
\end{minipage} \\
\midrule\noalign{}
\endhead
\bottomrule\noalign{}
\endlastfoot
\texttt{ICR\_00\_csv\_conformite\_report.csv} & CSV & Rapport de
conformité CSV (schéma + parsing) avant traitement. \\
\texttt{ICR\_00\_csv\_conformite\_problems.csv} & CSV & Détails des
problèmes de parsing (si présents). \\
\texttt{ICR\_donnees\_preparees.csv} & CSV & Données nettoyées et
standardisées \\
\texttt{ICR\_resume\_tous\_sites.csv} & CSV & Synthèse complétude +
représentativité tous sites \\
\texttt{ICR\_metrics\_pertinence\_tous\_sites.csv} & CSV & Métriques de
pertinence scientifique multi-sites \\
\texttt{ICR\_comparaison\_completude\_sites.png} & Graphique &
Comparaison du ratio de complétude entre sites \\
\texttt{ICR\_comparaison\_ir\_sites.png} & Graphique & Comparaison de
l'indice \(Ir\) entre sites \\
\texttt{ICR\_metrics\_pertinence\_heatmap\_sites.png} & Graphique &
Heatmap comparative des 7 métriques de pertinence \\
\texttt{ICR\_metrics\_pertinence\_score\_sites.png} & Graphique &
Classement global des sites par score de pertinence \\
\end{longtable}
}

\paragraph{\texorpdfstring{Fichiers par site
(\texttt{results/\textless{}nom\_du\_site\textgreater{}/})}{Fichiers par site (results/\textless nom\_du\_site\textgreater/)}}\label{fichiers-par-site-resultsnom_du_site}

{\def\LTcaptype{none} % do not increment counter
\begin{longtable}[]{@{}
  >{\raggedright\arraybackslash}p{(\linewidth - 6\tabcolsep) * \real{0.2500}}
  >{\raggedright\arraybackslash}p{(\linewidth - 6\tabcolsep) * \real{0.2500}}
  >{\raggedright\arraybackslash}p{(\linewidth - 6\tabcolsep) * \real{0.2500}}
  >{\raggedright\arraybackslash}p{(\linewidth - 6\tabcolsep) * \real{0.2500}}@{}}
\toprule\noalign{}
\begin{minipage}[b]{\linewidth}\raggedright
Fichier
\end{minipage} & \begin{minipage}[b]{\linewidth}\raggedright
Type
\end{minipage} & \begin{minipage}[b]{\linewidth}\raggedright
Condition
\end{minipage} & \begin{minipage}[b]{\linewidth}\raggedright
Description
\end{minipage} \\
\midrule\noalign{}
\endhead
\bottomrule\noalign{}
\endlastfoot
\texttt{ICR\_01\_courbe\_temps\_especes.csv} & CSV & Toujours & Richesse
observée et cumulée par visite \\
\texttt{ICR\_02\_tee\_ir.csv} & CSV & Toujours & Évolution de TEE et
\(Ir\) au fil des visites \\
\texttt{ICR\_03\_frequence\_especes.csv} & CSV & Toujours & Fréquences
d'occurrence des espèces \\
\texttt{ICR\_04\_resume\_site.csv} & CSV & Toujours & Résumé synthétique
du site \\
\texttt{ICR\_05\_modeles.csv} & CSV & Si modèles ajustés & Paramètres
des modèles linéaire/hyperbolique \\
\texttt{ICR\_06\_occupation\_spatiale.csv} & CSV & Si \texttt{placette}
disponible & Occupation spatiale par placette \\
\texttt{ICR\_07\_metrics\_pertinence.csv} & CSV & Toujours & 15
métriques de pertinence scientifique \\
\texttt{ICR\_01\_richesse\_par\_visite\_et\_cumul.png} & Graphique &
Toujours & Richesse par visite + richesse cumulée \\
\texttt{ICR\_02\_courbe\_temps\_especes\_hyperbole.png} & Graphique &
Toujours & Courbe temps-espèces + ajustement hyperbolique \\
\texttt{ICR\_03\_tee\_ir.png} & Graphique & Toujours & TEE et \(Ir\) en
deux panneaux \\
\texttt{ICR\_04\_histogramme\_frequences.png} & Graphique & Toujours &
Distribution des fréquences d'observabilité \\
\texttt{ICR\_05\_temporel\_vs\_spatial.png} & Graphique & Si
\texttt{placette} disponible & Fréquence temporelle vs occupation
spatiale \\
\texttt{ICR\_06\_CA\_placettes\_especes.png} & Graphique & Si
\texttt{vegan} + \texttt{placette} & Ordination CA/AFC
placettes-espèces \\
\texttt{ICR\_07\_metrics\_pertinence\_dashboard.png} & Graphique &
Toujours & Dashboard des 7 métriques de pertinence du site \\
\end{longtable}
}

\begin{quote}
\textbf{Convention de nommage :} tous les fichiers générés sont préfixés
\texttt{ICR\_} pour les distinguer des sorties des autres scripts du
projet.
\end{quote}

\begin{center}\rule{0.5\linewidth}{0.5pt}\end{center}

\subsubsection{Installation des
dépendances}\label{installation-des-duxe9pendances}

\paragraph{Packages requis}\label{packages-requis-1}

{\def\LTcaptype{none} % do not increment counter
\begin{longtable}[]{@{}ll@{}}
\toprule\noalign{}
Package & Rôle \\
\midrule\noalign{}
\endhead
\bottomrule\noalign{}
\endlastfoot
\texttt{dplyr} & Manipulation de données tabulaires \\
\texttt{tidyr} & Restructuration des données (format long/large) \\
\texttt{ggplot2} & Génération des graphiques \\
\texttt{purrr} & Itérations fonctionnelles sur les sites \\
\texttt{readr} & Lecture/écriture des fichiers CSV/TSV \\
\texttt{stringr} & Traitement des chaînes de caractères \\
\texttt{forcats} & Gestion des facteurs \\
\texttt{tibble} & Structures tabulaires modernes \\
\texttt{scales} & Mise en forme des axes graphiques \\
\texttt{gridExtra} & Assemblage de graphiques multi-panneaux \\
\end{longtable}
}

\begin{verbatim}
install.packages(c(
  "dplyr", "tidyr", "ggplot2", "purrr",
  "readr", "stringr", "forcats", "tibble",
  "scales", "gridExtra"
))
\end{verbatim}

\paragraph{Packages optionnels}\label{packages-optionnels}

{\def\LTcaptype{none} % do not increment counter
\begin{longtable}[]{@{}
  >{\raggedright\arraybackslash}p{(\linewidth - 2\tabcolsep) * \real{0.5000}}
  >{\raggedright\arraybackslash}p{(\linewidth - 2\tabcolsep) * \real{0.5000}}@{}}
\toprule\noalign{}
\begin{minipage}[b]{\linewidth}\raggedright
Package
\end{minipage} & \begin{minipage}[b]{\linewidth}\raggedright
Fonctionnalité activée
\end{minipage} \\
\midrule\noalign{}
\endhead
\bottomrule\noalign{}
\endlastfoot
\texttt{minpack.lm} & Ajustement hyperbolique (estimation d'asymptote,
ratio de complétude) \\
\texttt{vegan} & Ordination CA/AFC (analyse placettes--espèces) \\
\end{longtable}
}

\begin{verbatim}
install.packages(c("minpack.lm", "vegan"))
\end{verbatim}

\begin{quote}
Le script vérifie les packages requis au démarrage et affiche une
commande d'installation explicite si certains sont manquants.
L'exécution s'arrête proprement avec un message clair plutôt que de
produire des résultats partiels silencieux.
\end{quote}

\begin{center}\rule{0.5\linewidth}{0.5pt}\end{center}

\subsubsection{⚙️ Configuration}\label{configuration}

La configuration est centralisée dans l'objet \texttt{CONFIG} au début
du script
(\texttt{scripts/Inventaires\_completude\_representativite.R}). Toute
adaptation se fait directement dans ce bloc.

{\def\LTcaptype{none} % do not increment counter
\begin{longtable}[]{@{}
  >{\raggedright\arraybackslash}p{(\linewidth - 4\tabcolsep) * \real{0.3333}}
  >{\raggedright\arraybackslash}p{(\linewidth - 4\tabcolsep) * \real{0.3333}}
  >{\raggedright\arraybackslash}p{(\linewidth - 4\tabcolsep) * \real{0.3333}}@{}}
\toprule\noalign{}
\begin{minipage}[b]{\linewidth}\raggedright
Paramètre
\end{minipage} & \begin{minipage}[b]{\linewidth}\raggedright
Valeur par défaut
\end{minipage} & \begin{minipage}[b]{\linewidth}\raggedright
Description
\end{minipage} \\
\midrule\noalign{}
\endhead
\bottomrule\noalign{}
\endlastfoot
\texttt{input\_file} & \texttt{data/observations.csv} & Chemin du
fichier d'entrée (relatif à la racine du projet) \\
\texttt{output\_dir} & \texttt{results/ICR} & Dossier de sortie des CSV
et graphiques \\
\texttt{date\_format} & \texttt{\%Y-\%m-\%d} & Format de date R attendu
(ex. \texttt{2026-03-18}) \\
\texttt{csv\_strict\_mode} & \texttt{TRUE} & Stoppe l'exécution si le
CSV n'est pas conforme \\
\texttt{csv\_allow\_extra\_cols} & \texttt{FALSE} & Autorise/refuse les
colonnes supplémentaires \\
\texttt{csv\_required\_cols} &
\texttt{c("site","date","visite\_id","espece")} & Colonnes obligatoires
du schéma CSV \\
\texttt{csv\_optional\_cols} & \texttt{c("placette")} & Colonnes
optionnelles autorisées \\
\texttt{make\_ca} & \texttt{TRUE} & Active l'ordination CA/AFC si les
conditions sont réunies \\
\texttt{min\_visits\_for\_model} & \texttt{5} & Nombre minimal de
visites pour ajuster les modèles \\
\texttt{freq\_breaks} & \texttt{c(0,\ .10,\ .25,\ .50,\ .75,\ 1)} &
Bornes des classes de fréquence (en proportion) \\
\texttt{freq\_labels} & \texttt{exceptionnelle\ …\ constante} & Noms des
5 classes de fréquence \\
\texttt{width} & \texttt{10} & Largeur des graphiques exportés (en
pouces) \\
\texttt{height} & \texttt{6} & Hauteur des graphiques exportés (en
pouces) \\
\texttt{dpi} & \texttt{300} & Résolution des PNG exportés (points par
pouce) \\
\end{longtable}
}

\textbf{Exemple de personnalisation :}

\begin{verbatim}
# Dans le script, modifier CONFIG avant l'exécution :
CONFIG$input_file <- "data/mon_inventaire_2026.csv"
CONFIG$date_format <- "%d/%m/%Y"   # si vos dates sont au format jour/mois/année
CONFIG$make_ca <- FALSE            # désactiver la CA/AFC
CONFIG$min_visits_for_model <- 3   # assouplir le seuil si peu de visites
\end{verbatim}

\textbf{Variables d'environnement utiles :}

\begin{itemize}
\tightlist
\item
  \texttt{INVENTAIRES\_INPUT\_FILE}
\item
  \texttt{INVENTAIRES\_OUTPUT\_DIR}
\item
  \texttt{INVENTAIRES\_DATE\_FORMAT}
\item
  \texttt{INVENTAIRES\_CSV\_STRICT}
\item
  \texttt{INVENTAIRES\_CSV\_ALLOW\_EXTRA\_COLS}
\end{itemize}

\begin{center}\rule{0.5\linewidth}{0.5pt}\end{center}

\subsubsection{Interprétation des
indicateurs}\label{interpruxe9tation-des-indicateurs}

\paragraph{TEE -- Taux d'Espèces
Exceptionnelles}\label{tee-taux-despuxe8ces-exceptionnelles}

Proportion d'espèces observées \textbf{une seule fois} dans
l'inventaire.

{\def\LTcaptype{none} % do not increment counter
\begin{longtable}[]{@{}ll@{}}
\toprule\noalign{}
TEE & Interprétation \\
\midrule\noalign{}
\endhead
\bottomrule\noalign{}
\endlastfoot
\textgreater{} 0,40 & Inventaire souvent \textbf{incomplet} \\
0,20 -- 0,40 & Inventaire \textbf{intermédiaire} \\
\textless{} 0,20 & Inventaire généralement \textbf{bien consolidé} \\
\end{longtable}
}

\begin{itemize}
\tightlist
\item
  \textbf{TEE élevé en fin de suivi} : beaucoup d'espèces vues qu'une
  fois → inventaire en phase d'accumulation.
\item
  \textbf{TEE qui diminue} : signe positif, les espèces ``uniques''
  deviennent progressivement récurrentes.
\item
  \textbf{TEE stable et bas} : communauté mieux couverte, effort proche
  d'un plateau.
\end{itemize}

\paragraph{Ir -- Indice de
représentativité}\label{ir-indice-de-repruxe9sentativituxe9}

\[Ir = 1 - TEE\]

{\def\LTcaptype{none} % do not increment counter
\begin{longtable}[]{@{}ll@{}}
\toprule\noalign{}
\(Ir\) & Représentativité \\
\midrule\noalign{}
\endhead
\bottomrule\noalign{}
\endlastfoot
≥ 0,80 & Bonne à élevée \\
0,60 -- 0,80 & Moyenne \\
\textless{} 0,60 & Faible \\
\end{longtable}
}

\begin{itemize}
\tightlist
\item
  \textbf{Ir en hausse au fil des visites} : le protocole capture mieux
  la diversité réelle du site.
\item
  \textbf{Ir qui plafonne tôt} : effort potentiellement suffisant, sauf
  si période écologique incomplète.
\item
  \textbf{Ir faible et stagnant} : effort, calendrier ou stratification
  spatiale à revoir.
\end{itemize}

\paragraph{Complétude / asymptote
hyperbolique}\label{compluxe9tude-asymptote-hyperbolique}

\[\\text{Complétude} = \\frac{S\_{obs}}{S\_{asymptote}}\]

{\def\LTcaptype{none} % do not increment counter
\begin{longtable}[]{@{}ll@{}}
\toprule\noalign{}
Complétude & Interprétation \\
\midrule\noalign{}
\endhead
\bottomrule\noalign{}
\endlastfoot
\textgreater{} 0,90 & Inventaire très avancé \\
0,70 -- 0,90 & Progression correcte, inventaire partiel \\
\textless{} 0,70 & Insuffisante pour conclure sur l'exhaustivité \\
\end{longtable}
}

\begin{quote}
L'asymptote ne doit pas être interprétée comme une ``vérité absolue'',
mais comme un \textbf{repère quantitatif} de complétude.
\end{quote}

\paragraph{Score de pertinence scientifique (agrégé,
v1.4)}\label{score-de-pertinence-scientifique-agruxe9guxe9-v1.4}

\[\\text{Score} = 0{,}35 \\times Ir + 0{,}35 \\times \\text{complétude} + 0{,}20 \\times \\text{stabilité} + 0{,}10 \\times R^2\]

\begin{center}\rule{0.5\linewidth}{0.5pt}\end{center}

\subsubsection{📈 Lecture des graphiques}\label{lecture-des-graphiques}

\paragraph{\texorpdfstring{\texttt{ICR\_01\_richesse\_par\_visite\_et\_cumul.png}}{ICR\_01\_richesse\_par\_visite\_et\_cumul.png}}\label{icr_01_richesse_par_visite_et_cumul.png}

Combine \textbf{barres} (richesse observée par visite) et \textbf{ligne}
(richesse cumulée).

\begin{itemize}
\tightlist
\item
  \textbf{Barres encore hautes + cumul qui monte vite} → inventaire en
  phase active de découverte.
\item
  \textbf{Barres qui baissent + cumul qui se tasse} → entrée dans une
  phase de saturation.
\item
  \textbf{Dents de scie marquées} → forte dépendance météo / saison /
  observateur.
\item
  Signal d'alerte : si les dernières visites apportent encore beaucoup
  de nouvelles espèces, l'arrêt de l'inventaire peut être prématuré.
\end{itemize}

\paragraph{\texorpdfstring{\texttt{ICR\_02\_courbe\_temps\_especes\_hyperbole.png}}{ICR\_02\_courbe\_temps\_especes\_hyperbole.png}}\label{icr_02_courbe_temps_especes_hyperbole.png}

Superpose les observations à un ajustement hyperbolique (si
\texttt{minpack.lm} disponible) et une asymptote estimée.

\begin{itemize}
\tightlist
\item
  \textbf{Courbe observée proche de la courbe ajustée} → modèle
  cohérent, asymptote fiable.
\item
  \textbf{Écart important} → structure de données non compatible avec ce
  modèle (saisonnalité forte, ruptures d'effort, etc.).
\item
  \textbf{Asymptote nettement au-dessus de la richesse actuelle} → marge
  de découverte encore importante.
\end{itemize}

\paragraph{\texorpdfstring{\texttt{ICR\_03\_tee\_ir.png}}{ICR\_03\_tee\_ir.png}}\label{icr_03_tee_ir.png}

Deux panneaux complémentaires :

panneau haut : nombre d'espèces observées une seule fois
(\texttt{nb\_esp\_1fois})

panneau bas : indice \(Ir\) sur l'échelle \(\[0,1\]\) avec seuils
visuels (0,60 et 0,80)

\textbf{Espèces uniques qui ralentissent} → signal de consolidation de
l'inventaire.

\textbf{Ir qui augmente puis se stabilise} → représentativité qui
converge.

Les seuils sont des repères empiriques : toujours interpréter avec le
contexte (saison, effort, habitat).

\paragraph{\texorpdfstring{\texttt{ICR\_04\_histogramme\_frequences.png}}{ICR\_04\_histogramme\_frequences.png}}\label{icr_04_histogramme_frequences.png}

Distribution du nombre d'espèces selon leur fréquence d'occurrence.

\begin{itemize}
\tightlist
\item
  \textbf{Pic sur les faibles fréquences} → communauté riche en espèces
  rares/occasionnelles.
\item
  \textbf{Répartition plus équilibrée} → noyau d'espèces régulières bien
  documenté.
\item
  \textbf{Queue vers les fortes fréquences} → présence d'espèces
  constantes, potentiels bons indicateurs écologiques.
\end{itemize}

\paragraph{\texorpdfstring{\texttt{ICR\_05\_temporel\_vs\_spatial.png}
\emph{(si} \texttt{\_placette\_}
\emph{disponible)}}{ICR\_05\_temporel\_vs\_spatial.png (si \_placette\_ disponible)}}\label{icr_05_temporel_vs_spatial.png-si-_placette_-disponible}

Nuage de points reliant fréquence temporelle (\(x\) = proportion de
visites) et occupation spatiale (\(y\) = proportion de placettes).

\begin{itemize}
\tightlist
\item
  \textbf{Haut-droite} → espèces fréquentes et largement réparties
  (généralistes / communes).
\item
  \textbf{Bas-gauche} → espèces rares et localisées (vraie rareté).
\item
  \textbf{Bas-droite} → espèces temporellement fréquentes mais
  spatialement confinées (micro-habitats stables).
\item
  \textbf{Haut-gauche} → espèces spatialement larges mais irrégulières
  dans le temps (phénologie / pulses).
\end{itemize}

Ce diagramme permet de distinguer rareté vraie, rareté d'observation et
hétérogénéité spatiale.

\paragraph{\texorpdfstring{\texttt{ICR\_06\_CA\_placettes\_especes.png}
\emph{(si} \texttt{\_vegan\_} \emph{+} \texttt{\_placette\_}
\emph{disponibles)}}{ICR\_06\_CA\_placettes\_especes.png (si \_vegan\_ + \_placette\_ disponibles)}}\label{icr_06_ca_placettes_especes.png-si-_vegan_-_placette_-disponibles}

Ordination CA/AFC explorant les associations placettes--espèces.

\begin{itemize}
\tightlist
\item
  \textbf{Points proches} → profils mycologiques similaires.
\item
  \textbf{Groupes compacts} → unités écologiques homogènes.
\item
  \textbf{Points isolés} → placettes atypiques ou espèces spécialisées.
\item
  Les axes de CA sont relatifs au jeu de données analysé ; comparer des
  CA entre jeux différents demande prudence.
\end{itemize}

\paragraph{\texorpdfstring{\texttt{ICR\_07\_metrics\_pertinence\_dashboard.png}}{ICR\_07\_metrics\_pertinence\_dashboard.png}}\label{icr_07_metrics_pertinence_dashboard.png}

Dashboard panoramique par site : barres horizontales pour chacun des 7
indicateurs (normalisés 0--1).

\begin{itemize}
\tightlist
\item
  Lignes de seuil pointillées à \textbf{0,60} (acceptable) et
  \textbf{0,80} (bon).
\item
  Tout indicateur sous 0,60 motive un complément d'effort ou une
  révision du protocole.
\end{itemize}

\paragraph{Graphiques multi-sites}\label{graphiques-multi-sites}

{\def\LTcaptype{none} % do not increment counter
\begin{longtable}[]{@{}
  >{\raggedright\arraybackslash}p{(\linewidth - 2\tabcolsep) * \real{0.5000}}
  >{\raggedright\arraybackslash}p{(\linewidth - 2\tabcolsep) * \real{0.5000}}@{}}
\toprule\noalign{}
\begin{minipage}[b]{\linewidth}\raggedright
Graphique
\end{minipage} & \begin{minipage}[b]{\linewidth}\raggedright
Ce qu'il montre
\end{minipage} \\
\midrule\noalign{}
\endhead
\bottomrule\noalign{}
\endlastfoot
\texttt{ICR\_comparaison\_completude\_sites.png} & Sites classés par
ratio de complétude --- identifier rapidement les sites à renforcer \\
\texttt{ICR\_comparaison\_ir\_sites.png} & Sites classés par \(Ir\)
final --- prioriser les revisites (sites à \(Ir\) faible en premier) \\
\texttt{ICR\_metrics\_pertinence\_heatmap\_sites.png} & Heatmap 7
métriques × tous sites (rouge = faible, vert = bon) --- vue synthétique
comparative \\
\texttt{ICR\_metrics\_pertinence\_score\_sites.png} & Classement par
score agrégé avec zones rouge/jaune/vert --- synthèse exécutive \\
\end{longtable}
}

\begin{center}\rule{0.5\linewidth}{0.5pt}\end{center}

\subsubsection{Bonnes pratiques
d'interprétation}\label{bonnes-pratiques-dinterpruxe9tation}

\begin{itemize}
\tightlist
\item
  Interpréter \textbf{ensemble} : courbe cumulée + TEE/Ir + fréquences +
  spatial.
\item
  Toujours contextualiser avec : saison, pression d'échantillonnage,
  météo, observateurs.
\item
  Préférer les \textbf{tendances sur plusieurs visites} plutôt qu'une
  lecture ponctuelle.
\item
  Un site bas en complétude n'est pas ``mauvais'' : il peut être plus
  hétérogène ou plus difficile à prospecter.
\item
  Croiser complétude et \(Ir\) avant de conclure sur l'effort à fournir
  : un \(Ir\) élevé avec une complétude faible peut signifier un site
  très riche, pas un protocole défaillant.
\item
  Pour des décisions de gestion, combiner ces indicateurs avec
  l'expertise terrain et la connaissance taxonomique.
\end{itemize}

\begin{center}\rule{0.5\linewidth}{0.5pt}\end{center}

\subsubsection{Dépannage (FAQ)}\label{duxe9pannage-faq}

\textbf{Le script ne trouve pas mon fichier d'entrée}\\
Vérifier qu'un fichier existe dans \texttt{data/} sous :
\texttt{observations.csv}, \texttt{observations.txt} ou
\texttt{observations.tsv}, et que les 4 colonnes obligatoires
(\texttt{site}, \texttt{date}, \texttt{visite\_id}, \texttt{espece})
sont présentes.

\textbf{Le script s'arrête avec ``CSV non conforme''}\\
Consulter :

\begin{itemize}
\tightlist
\item
  \texttt{results/ICR/ICR\_00\_csv\_conformite\_report.csv}
\item
  \texttt{results/ICR/ICR\_00\_csv\_conformite\_problems.csv} (si
  présent)
\end{itemize}

Causes typiques : colonne obligatoire manquante, colonne supplémentaire
non autorisée, ou ligne mal formée.

\textbf{Erreur ``Packages requis manquants''}\\
Le script affiche la commande \texttt{install.packages(...)} exacte à
exécuter. Copier-coller cette commande dans R, puis relancer le script.

\textbf{Erreur sur les dates}\\
Contrôler le format de la colonne \texttt{date} et adapter
\texttt{CONFIG\$date\_format} dans le script (défaut :
\texttt{\%Y-\%m-\%d}, ex. \texttt{2026-03-23}). Format courant
alternatif : \texttt{\%d/\%m/\%Y} pour \texttt{23/03/2026}.

\textbf{Le modèle hyperbolique n'est pas calculé}\\
Installer \texttt{minpack.lm}. Sans ce package, l'analyse continue mais
sans asymptote hyperbolique ni ratio de complétude. Vérifier aussi que
le site a au moins \texttt{CONFIG\$min\_visits\_for\_model} visites
(défaut : 5).

\textbf{La CA/AFC n'apparaît pas}\\
Nécessite : package \texttt{vegan} installé + colonne \texttt{placette}
présente + matrice placette-espèce suffisamment informative (plusieurs
placettes avec plusieurs espèces différentes). Vérifier également que
\texttt{CONFIG\$make\_ca} est \texttt{TRUE}.

\textbf{Le manifeste signale des fichiers MANQUANT}

\begin{itemize}
\tightlist
\item
  \texttt{MANQUANT} : fichier obligatoire non produit --- voir les logs
  d'exécution pour l'erreur associée.
\item
  \texttt{OPTIONNEL\_NON\_GENERE} : normal si les conditions ne sont pas
  réunies (ex. \texttt{placette} absent, \texttt{vegan} non installé,
  \texttt{make\_ca\ =\ FALSE}).
\end{itemize}

\textbf{Les graphiques de métriques de pertinence ne sont pas générés}\\
Le site doit avoir au moins 2 visites et les colonnes \texttt{Ir\_final}
et \texttt{completude} doivent être présentes dans le résumé du site.

\textbf{Exécution très lente}\\
Activer \texttt{BENCHMARK\_MODE\ \textless{}-\ TRUE} en tête du script
pour identifier les goulots d'étranglement. Sur de grands jeux de
données, la phase de calculs cumulés bénéficie de la vectorisation
introduite en v1.2.

\begin{center}\rule{0.5\linewidth}{0.5pt}\end{center}

\subsubsection{Historique des versions}\label{historique-des-versions}

{\def\LTcaptype{none} % do not increment counter
\begin{longtable}[]{@{}
  >{\raggedright\arraybackslash}p{(\linewidth - 4\tabcolsep) * \real{0.3333}}
  >{\raggedright\arraybackslash}p{(\linewidth - 4\tabcolsep) * \real{0.3333}}
  >{\raggedright\arraybackslash}p{(\linewidth - 4\tabcolsep) * \real{0.3333}}@{}}
\toprule\noalign{}
\begin{minipage}[b]{\linewidth}\raggedright
Version
\end{minipage} & \begin{minipage}[b]{\linewidth}\raggedright
Phase
\end{minipage} & \begin{minipage}[b]{\linewidth}\raggedright
Principales évolutions
\end{minipage} \\
\midrule\noalign{}
\endhead
\bottomrule\noalign{}
\endlastfoot
\textbf{v1.4} (actuelle) & Phase 4 --- Industrialisation & Métriques de
pertinence scientifique (stabilité, robustesse, cohérence, score
agrégé), rapport Quarto, manifeste automatique des sorties \\
\textbf{v1.3} & Phase 3 --- Graphiques & Thème unifié CVD-friendly
(palette accessible 5 couleurs), annotations scientifiques, diagramme
TEE/Ir 2 panneaux avec seuils, comparaisons sites annotées, export
multi-formats (PNG, PDF, SVG) \\
\textbf{v1.2} & Phase 2 --- Performance & Vectorisation des calculs
cumulés, mode benchmarking (\texttt{BENCHMARK\_MODE}) \\
\textbf{v1.1} & Phase 1 --- Fiabilisation & Validation robuste des
données d'entrée, détection automatique du séparateur (CSV/TSV/TXT),
messages d'erreur explicites \\
\end{longtable}
}

\begin{center}\rule{0.5\linewidth}{0.5pt}\end{center}

\subsubsection{Glossaire}\label{glossaire}

{\def\LTcaptype{none} % do not increment counter
\begin{longtable}[]{@{}
  >{\raggedright\arraybackslash}p{(\linewidth - 2\tabcolsep) * \real{0.5000}}
  >{\raggedright\arraybackslash}p{(\linewidth - 2\tabcolsep) * \real{0.5000}}@{}}
\toprule\noalign{}
\begin{minipage}[b]{\linewidth}\raggedright
Terme
\end{minipage} & \begin{minipage}[b]{\linewidth}\raggedright
Définition
\end{minipage} \\
\midrule\noalign{}
\endhead
\bottomrule\noalign{}
\endlastfoot
\textbf{TEE} & Taux d'Espèces Exceptionnelles : proportion d'espèces
observées une seule fois \\
\textbf{Ir} & Indice de représentativité : \(Ir = 1 - TEE\) \\
\textbf{Asymptote} & Richesse théorique maximale estimée par le modèle
hyperbolique \\
\textbf{Complétude} & Ratio \(S\_{obs} / S\_{asymptote}\) : avancement
de l'inventaire vers la richesse théorique \\
\textbf{Robustesse} & Qualité d'ajustement (\(R^2\)) du modèle
d'asymptote \\
\textbf{Stabilité} & Convergence de \(Ir\) sur les dernières visites \\
\textbf{Cohérence} & Corrélation Spearman entre fréquences temporelles
et spatiales, normalisée en {[}0,1{]} \\
\textbf{Score pertinence} & Indicateur agrégé pondéré :
\(0{,}35 \\times Ir + 0{,}35 \\times \\text{complétude} + 0{,}20 \\times \\text{stabilité} + 0{,}10 \\times R^2\) \\
\textbf{CA / AFC} & Correspondence Analysis / Analyse Factorielle des
Correspondances \\
\textbf{CSV} & Comma-Separated Values : fichier texte tabulaire, champs
séparés par des virgules \\
\textbf{TSV} & Tab-Separated Values : fichier texte tabulaire, champs
séparés par des tabulations \\
\textbf{ICR} & Préfixe de nommage de tous les fichiers produits par ce
script \\
\textbf{CVD-friendly} & Compatible daltoniens (Color Vision Deficiency)
: palette graphique accessible \\
\textbf{vegan} & Package R (\emph{Vegetation Ecology And other Nature
data Group Analysis}) : outils d'analyse en écologie des communautés ---
ordinations (CA/AFC, PCA, NMDS), indices de diversité, tests de
permutation, etc. Dans ce projet, utilisé pour l'ordination CA/AFC entre
placettes et espèces \\
\end{longtable}
}

\begin{center}\rule{0.5\linewidth}{0.5pt}\end{center}

\subsection{Références}\label{ruxe9fuxe9rences}

\begin{itemize}
\tightlist
\item
  Cours 24 Inventaires (diapositives N° 25-48 du ) - DU Mycologie 2026
  du Professeur Pierre-Arthur Moreau, Université de Lille (UFR3S PHAR)
\item
  Méthodes et définitions utilisées dans ce document (TEE, complétude,
  CA/AFC, cohérence tempo-spatiale) s'appuient sur les principes
  standards de l'écologie des communautés et de l'analyse multivariée
  sous R, en particulier via les packages \textbf{vegan}, \textbf{dplyr}
  et \textbf{ggplot2}. Pour un cadre théorique détaillé, se référer à la
  documentation officielle de ces packages ainsi qu'aux ouvrages de
  référence en ordination écologique.
\end{itemize}

\begin{center}\rule{0.5\linewidth}{0.5pt}\end{center}

\subsection{Contact}\label{contact}

\textbf{Auteur :} Eddy Boite\\
Pour toute question ou évolution, ouvrir une issue sur ce dépôt.
